\documentclass[10pt,a4paper]{amsart}
\usepackage[margin=19mm]{geometry}
\usepackage{amsmath,amssymb,mathtools}
\usepackage[scr=rsfs,cal=euler]{mathalfa}
\usepackage{xfrac}
\usepackage[unicode,hidelinks,bookmarks=false]{hyperref}
\newcommand{\ie}{i.e.\ }
\newcommand{\cf}{cf.\ }
\newcommand{\resp}{respectively\ }
\newcommand{\Subsection}{\subsection}
\providecommand{\nobreakdash}{\nobreak\hbox{-}}
\setlength{\emergencystretch}{2em}
\multlinegap0pt
\usepackage{bm}
\usepackage{bbm}
\usepackage{dsfont}
\usepackage{xspace}
\usepackage{cancel}
\usepackage{mathtools}
\usepackage{amsthm}
\makeatletter
\@ifundefined{theorem}{\newtheorem{theorem}{Theorem}}{}
\@ifundefined{lemma}{\newtheorem{lemma}[theorem]{Lemma}}{}
\@ifundefined{proposition}{\newtheorem{proposition}[theorem]{Proposition}}{}
\makeatother
\newcommand{\Ima}{\operatorname{Im}}
\newcommand{\M}{\mathsf M}
\newcommand{\Rea}{\operatorname{Re}}
\newcommand{\R}{\mathbb R}
\newcommand{\dd}{\,\mathrm d}
\title[Finite-mass soliton-type rigidity and four-channel reduction]{Finite-mass soliton-type rigidity and four-channel reduction for the three-dimensional nonradial focusing energy-critical nonlinear Schr\"odinger equation}
\author{Pang-Hung Chung, Dan Han}
\date{Source: arXiv:2607.25735v1 (CC BY 4.0)}
\begin{document}
\maketitle

\noindent\emph{Source and licence.} Finite-mass soliton-type rigidity and four-channel reduction for the three-dimensional nonradial focusing energy-critical nonlinear Schr\"odinger equation. Version arXiv:2607.25735v1 (CC BY 4.0), distributed under the Creative Commons Attribution 4.0 International licence, as recorded in the arXiv licence metadata for this exact version and shown on the abstract page https://arxiv.org/abs/2607.25735v1. Full rights notice: licenses/arxiv-2607.25735-CC-BY-4.0.txt.\par\smallskip

\noindent\textit{Editorial scope.} Counted unit: the Proposition whose source label is \texttt{prop:exact-Morawetz} in the article (source0.tex lines 2044--2067, source label \texttt{prop:exact-Morawetz}), delivered together with the article's own complete original proof of it (source0.tex lines 2069--2154, one \texttt{proof} environment of 86 lines). No mathematics is written, repaired or re-derived: statement and proof are the article's own text line for line. Required context retained uncounted: lem:kernel-bounds (lines 1868--1890); eq:A-endpoint-bound (lines 1879--1889); eq:T-def (lines 1943--1949); lem:local-laws (lines 1959--1969); eq:local-mass-law (lines 1961--1963); eq:local-momentum-law (lines 1965--1968); eq:Morawetz-functional (lines 2036--2040). Every definition, display and label of those lines is retained unchanged. Omitted from this package and not redistributed: 1--1867, 1890--1942, 1950--1958, 1964--1964, 1969--2035, 2041--2043, 2068--2068, 2155--3070. Nothing inside a retained excerpt is omitted or altered: every retained body is delivered exactly as the article's lines are written, so no omission receipt of any kind accompanies this package. Citations: the retained excerpts carry no citation command at all, so no bibliography entry of the article is transcribed and no reference is redistributed. Reference adaptation: none was needed and none was made; every reference inside the delivered unit resolves inside it, so no unresolved reference or citation is produced. Printed slips kept literal: no printed word, symbol, hypothesis or number of the counted statement or of its proof was altered. Layout and class: the article's own class is \texttt{interact} with packages including mathtools, enumitem, booktabs, tabularx, array, ragged2e, tikz, microtype, natbib, hyperref; the wrapper is the lane's pinned \texttt{amsart} editorial wrapper at 10pt on A4, which restates the theorem-like environments the retained text needs and adds the article's own macro definitions that the retained text needs (\texttt{\textbackslash Ima}, \texttt{\textbackslash M}, \texttt{\textbackslash R}, \texttt{\textbackslash Rea}, \texttt{\textbackslash dd}), with extra wrapper packages (bm, bbm, dsfont, xspace, cancel, mathtools). The delivered numbering is the unit's own. Assets: no retained span contains a figure environment, a graphics-inclusion command, a PDF-page inclusion, a table or a TikZ picture; no image file is copied into this package, the package declares no asset dependency and no image file is redistributed with it. Licence: version arXiv:2607.25735v1 of this article is distributed under the Creative Commons Attribution 4.0 International licence, as recorded in the arXiv licence metadata for this exact version and shown on the abstract page https://arxiv.org/abs/2607.25735v1. A full rights notice is delivered at licenses/arxiv-2607.25735-CC-BY-4.0.txt. Characters: every retained excerpt is pure ASCII, so the delivered engine, XeLaTeX with its default Latin Modern text font, reports no missing character.
\begin{lemma}[Kernel estimates]\label{lem:kernel-bounds}
For $z\in\R^3$,
\begin{align}
 \partial_kA_j(z)
 &=\delta_{jk}\phi(z)+z_j\partial_k\phi(z),
 \label{eq:A-Hessian}\\
 \operatorname{div}A(z)
 &=3\phi(z)+z\cdot\nabla\phi(z).
 \label{eq:A-divergence}
\end{align}
There exists $C$, depending only on $\zeta$, such that
\begin{align}
 \|A\|_\infty&\le Ce^JR_0,
 \label{eq:A-endpoint-bound}\\
 \|z\otimes\nabla\phi\|_\infty
 +\|z\cdot\nabla\phi\|_\infty
 &\le \frac{C}{J},
 \label{eq:flat-error-bound}\\
 \|\Delta\operatorname{div}A\|_\infty
 &\le \frac{C}{JR_0^2}.
 \label{eq:biharmonic-bound}
\end{align}
\end{lemma}

\begin{align}
 m[v](t,x)&:=|v(t,x)|^2,\notag\\
 p_j[v](t,x)&:=\Ima\bigl(\overline{v(t,x)}\,\partial_jv(t,x)\bigr),\notag\\
 T_{jk}[v](t,x)&:=\frac12\delta_{jk}\Delta m[v](t,x)
 -2\Rea\bigl(\partial_jv(t,x)\,\partial_k\overline{v(t,x)}\bigr).
 \label{eq:T-def}
\end{align}

\begin{lemma}[Local conservation laws]\label{lem:local-laws}
If $u$ is a smooth rapidly decaying solution, then
\begin{equation}\label{eq:local-mass-law}
 \partial_tm+2\partial_kp_k=0,
\end{equation}
and
\begin{equation}\label{eq:local-momentum-law}
 \partial_tp_j
 =\partial_kT_{jk}+\frac23\partial_j|u|^6.
\end{equation}
\end{lemma}

\begin{equation}\label{eq:Morawetz-functional}
 \mathcal I_{J,R_0}(t)
 :=2\iint_{\R^3\times\R^3}
 m(t,y)A_j(x-y)p_j(t,x)\dd x\dd y.
\end{equation}

\begin{proposition}[Morawetz identity]\label{prop:exact-Morawetz}
If $u$ is smooth and rapidly decaying, then
\begin{align}
 \frac{\dd}{\dd t}\mathcal I_{J,R_0}(t)
 ={}&4\iint \partial_kA_j(x-y)
 \Bigl[
 m(y)\Rea(\partial_ju\,\partial_k\bar u)(x)
 -p_j(x)p_k(y)
 \Bigr]\dd x\dd y
 \notag\\
 &-\frac43\iint
 (\operatorname{div}A)(x-y)m(y)|u(x)|^6\dd x\dd y
 \notag\\
 &-\iint
 \Delta(\operatorname{div}A)(x-y)m(x)m(y)\dd x\dd y.
 \label{eq:exact-Morawetz-identity}
\end{align}
Moreover, if $\M(u)<\infty$, then
\begin{equation}\label{eq:Morawetz-endpoint}
 \sup_t|\mathcal I_{J,R_0}(t)|
 \le Ce^JR_0\M(u)^{3/2}
 \sup_t\|\nabla u(t)\|_2.
\end{equation}
\end{proposition}

\begin{proof}
We suppress the time variable. Differentiating \eqref{eq:Morawetz-functional} and using the conservation laws in Lemma~\ref{lem:local-laws}, define
\begin{align*}
 I_m&:=2\iint \partial_tm(y)A_j(x-y)p_j(x)\dd x\dd y,\\
 I_p&:=2\iint m(y)A_j(x-y)\partial_tp_j(x)\dd x\dd y.
\end{align*}
Then
\begin{equation}\label{eq:I-prime-split}
 \mathcal I_{J,R_0}'=I_m+I_p.
\end{equation}
By \eqref{eq:local-mass-law}, integration by parts in $y$, together with
\[
 \partial_{y_k}A_j(x-y)=-\partial_kA_j(x-y),
\]
gives
\begin{align}
 I_m
 &=-4\iint \partial_{y_k}p_k(y)A_j(x-y)p_j(x)\dd x\dd y
 \notag\\
 &=4\iint p_k(y)\partial_{y_k}A_j(x-y)p_j(x)\dd x\dd y
 \notag\\
 &=-4\iint \partial_kA_j(x-y)p_j(x)p_k(y)\dd x\dd y.
 \label{eq:Im-computed}
\end{align}
By \eqref{eq:local-momentum-law},
\begin{align*}
 I_p
 &=2\iint m(y)A_j(x-y)\partial_kT_{jk}(x)\dd x\dd y\\
 &\quad+\frac43\iint m(y)A_j(x-y)\partial_j|u(x)|^6\dd x\dd y.
\end{align*}
Integrating by parts in $x$,
\begin{align}
 I_p
 &=-2\iint m(y)\partial_kA_j(x-y)T_{jk}(x)\dd x\dd y
 \notag\\
 &\quad-\frac43\iint m(y)(\operatorname{div}A)(x-y)|u(x)|^6\dd x\dd y.
 \label{eq:Ip-after-parts}
\end{align}
Using the stress tensor \eqref{eq:T-def},
\begin{align}
 -2\partial_kA_jT_{jk}
 &=-(\operatorname{div}A)\Delta m
 +4\partial_kA_j\Rea(\partial_ju\,\partial_k\bar u).
 \label{eq:stress-expanded}
\end{align}
For the first term in \eqref{eq:stress-expanded}, two further integrations by parts in $x$ yield
\begin{align}
 -\int_{\R^3}(\operatorname{div}A)(x-y)\Delta m(x)\dd x
 &=-\int_{\R^3}\Delta(\operatorname{div}A)(x-y)m(x)\dd x.
 \label{eq:double-parts}
\end{align}
Substituting \eqref{eq:stress-expanded} and \eqref{eq:double-parts} into \eqref{eq:Ip-after-parts} gives
\begin{align}
 I_p
 ={}&4\iint m(y)\partial_kA_j(x-y)
 \Rea(\partial_ju\,\partial_k\bar u)(x)\dd x\dd y
 \notag\\
 &-\frac43\iint
 (\operatorname{div}A)(x-y)m(y)|u(x)|^6\dd x\dd y
 \notag\\
 &-\iint
 \Delta(\operatorname{div}A)(x-y)m(x)m(y)\dd x\dd y.
 \label{eq:Ip-computed}
\end{align}
Substitution of \eqref{eq:Im-computed} and \eqref{eq:Ip-computed} into \eqref{eq:I-prime-split} proves \eqref{eq:exact-Morawetz-identity}.

For the endpoint bound, the pointwise estimate
\[
 |p(t,x)|\le |u(t,x)|\,|\nabla u(t,x)|
\]
and Cauchy--Schwarz give
\begin{equation}\label{eq:p-L1}
 \|p(t)\|_1
 \le\|u(t)\|_2\|\nabla u(t)\|_2
 =\M(u)^{1/2}\|\nabla u(t)\|_2.
\end{equation}
By Lemma~\ref{lem:kernel-bounds}, the endpoint estimate \eqref{eq:A-endpoint-bound} holds. Combining it with \eqref{eq:Morawetz-functional} and \eqref{eq:p-L1},
\begin{align*}
 |\mathcal I_{J,R_0}(t)|
 &\le2\|A\|_\infty
 \left(\int m(y)\dd y\right)
 \left(\int|p(x)|\dd x\right)\\
 &\le Ce^JR_0\M(u)^{3/2}\|\nabla u(t)\|_2.
\end{align*}
Taking the supremum over time proves \eqref{eq:Morawetz-endpoint}.
\end{proof}

\end{document}
